\documentclass[11pt,a4paper]{report}

% Packages
\usepackage[margin=0.85in]{geometry}
\usepackage{fancyhdr}
\usepackage{titlesec}
\usepackage{booktabs}
\usepackage{longtable}
\usepackage{array}
\usepackage{amsmath}
\usepackage{enumitem}
\usepackage{hyperref}
\usepackage{xcolor}
\usepackage{graphicx}
\usepackage{float}
\usepackage{verbatim}
\usepackage{listings}
\usepackage{tikz}
\usetikzlibrary{shapes.geometric, arrows.meta, positioning, calc, fit, backgrounds}

% Colors
\definecolor{urteg}{RGB}{0,128,128}
\definecolor{posixblue}{RGB}{0,51,102}
\definecolor{rustorange}{RGB}{183,65,14}
\definecolor{capellagreen}{RGB}{0,128,0}
\definecolor{tensorflow}{RGB}{255,140,0}

% Listings style
\lstdefinestyle{kernelschema}{
    basicstyle=\ttfamily\footnotesize,
    breaklines=true,
    frame=single,
    backgroundcolor=\color{gray!6},
    keywordstyle=\color{posixblue}\bfseries,
    commentstyle=\color{gray}
}

% Header/Footer
\pagestyle{fancy}
\fancyhf{}
\fancyhead[L]{\small URTE Kernel Threads - Hybrid Parallel Spec | Confidential}
\fancyhead[R]{\small June 2026 | v1.0}
\fancyfoot[C]{\thepage}
\renewcommand{\headrulewidth}{0.4pt}

% Section formatting
\titleformat{\chapter}{\Large\bfseries\color{posixblue}}{\thechapter.}{0.5em}{}
\titleformat{\section}{\large\bfseries\color{urteg}}{\thesection}{0.5em}{}
\titleformat{\subsection}{\normalsize\bfseries\color{rustorange}}{\thesubsection}{0.5em}{}

% Title
\title{\textbf{URTE Kernel Threads Environment}\\[0.3cm]
\large Hybrid TensorFlow + Vector Parallel Computing\\[0.2cm]
\normalsize for System Calls Add-in Implementation\\[0.2cm]
with \textbf{Capella} MBSE and \textbf{Rust} Compiler\\[0.3cm]
\textbf{Multispectral Parallel Thread Management System}}
\author{
    \textbf{Prepared by:}\\[0.1cm]
    Grok xAI Research \& UAB Saptera\\[0.1cm]
    Kristupas Gelzinis \& Olek Suchodolski
}
\date{June 2026}

\begin{document}

\maketitle

\begin{abstract}
This document details the \textbf{URTE Kernel Threads Environment} using a \textbf{hybrid TensorFlow + Vector Parallel Computing} model for efficient execution of POSIX and URTE-specific system calls.

The architecture is based on a \textbf{Multispectral Parallel Thread Management System} that operates across all scales (molecular to planetary) while maintaining full POSIX compliance. Implementation guidance is provided for:
\begin{itemize}
    \item \textbf{Rust} as the primary language for safe, high-performance kernel and user-space components
    \item \textbf{Capella} (Arcadia method) for formal MBSE modeling of the thread environment
\end{itemize}

All designs are presented with detailed \textbf{ASCII schemas} and are aligned with the previous URTE POSIX Kernel and Arcadia/Capella specifications.
\end{abstract}

\tableofcontents
\newpage

%==============================================================================
\chapter{Introduction to URTE Kernel Threads Environment}
%==============================================================================

\section{Overview}

The URTE Kernel requires a sophisticated threading model to handle:
\begin{itemize}
    \item Concurrent multi-scale digital twin simulations (PINNs across scales)
    \item Real-time intervention orchestration and TRL Pull acceleration
    \item Parallel execution of guardrail checks and ethical evaluations
    \item High-throughput telemetry ingestion from clinical, ecological, and space sensors
\end{itemize}

Traditional POSIX threading (pthreads) is extended with \textbf{hybrid parallel computing} capabilities:
\begin{itemize}
    \item \textbf{TensorFlow integration} for AI/ML workloads (generative design, reinforcement learning optimizers, PINN inference)
    \item \textbf{Vector/SIMD + GPU acceleration} for linear algebra operations on state vectors $u_s$
    \item \textbf{Multispectral thread pools} that dynamically allocate threads across scale levels
\end{itemize}

\section{Design Goals}

\begin{enumerate}
    \item Maintain strict \textbf{POSIX compliance} for all standard threading interfaces
    \item Provide \textbf{zero-cost abstractions} for URTE-specific parallel operations (Rust)
    \item Enable \textbf{formal modeling} in Capella using Arcadia viewpoints
    \item Support \textbf{hybrid execution}: CPU vector units + TensorFlow runtime + optional GPU/TPU offload
    \item Ensure \textbf{deterministic real-time behavior} for safety-critical paths while allowing best-effort parallel learning workloads
\end{enumerate}

%==============================================================================
\chapter{Multispectral Parallel Thread Management System}
%==============================================================================

\section{Thread Classification}

\begin{lstlisting}[style=kernelschema, caption={URTE Multispectral Thread Types}]
Thread Type                  | Scale Range     | Priority | Parallelism Model          | Language
-----------------------------|-----------------|----------|----------------------------|----------
POSIX Standard Thread        | N/A             | Normal   | 1:1 or M:N (Rust)          | Rust/C
Regen Agent Thread           | Any             | High     | Scale-aware pool           | Rust
Digital Twin Simulator       | Multi-scale     | Variable | TensorFlow + Vector SIMD   | Rust + TF
Guardrail Evaluator          | All             | Critical | Vector + Small TF models   | Rust
TRL Pull Coordinator         | Cross-scale     | High     | Event-driven + Parallel    | Rust
Telemetry Ingestor           | Any             | Real-time| Lock-free queues + Vector  | Rust
\end{lstlisting}

\section{Multispectral Thread Pool Architecture (ASCII Schema)}

\begin{lstlisting}[style=kernelschema, caption={Multispectral Thread Pool Manager}]
+-----------------------------------------------------------------------+
|                    URTE KERNEL THREAD MANAGER                         |
|-----------------------------------------------------------------------|
|  Global Thread Registry (POSIX compliant + URTE extensions)           |
|                                                                       |
|  +------------------+  +------------------+  +------------------+     |
|  | Molecular Pool   |  | Tissue/Organ     |  | Ecosystem/       |     |
|  | (Scale 0-1)      |  | Pool (Scale 2-3) |  | Planetary Pool   |     |
|  | Threads: 4-16    |  | Threads: 8-32    |  | (Scale 4-6)      |     |
|  | Vector Width:    |  | Vector Width:    |  | Threads: 2-8     |     |
|  | AVX-512 / NEON   |  | AVX-512 + TF     |  | TF Heavy         |     |
|  +------------------+  +------------------+  +------------------+     |
|           |                     |                     |               |
|           v                     v                     v               |
|  +----------------------------------------------------------------+   |
|  |           Hybrid Execution Engine (Rust + TensorFlow C API)    |   |
|  |  - TensorFlow Lite / Full TF for PINN & Generative models      |   |
|  |  - Rust rayon / crossbeam for CPU vector parallelism           |   |
|  |  - GPU offload via CUDA / Vulkan / Metal when available        |   |
|  +----------------------------------------------------------------+   |
|                                                                       |
|  Scheduler Integration: TRL Pull aware + POSIX SCHED_FIFO/RR/OTHER   |
+-----------------------------------------------------------------------+
\end{lstlisting}

%==============================================================================
\chapter{Hybrid TensorFlow + Vector Parallel Computing for System Calls}
%==============================================================================

\section{System Call Parallel Execution Model}

URTE-specific system calls (especially `urte_guardrail_check`, `urte_state_vector_create`, and `urte_trl_pull_notify`) are executed using a hybrid parallel pipeline:

\begin{lstlisting}[style=kernelschema, caption={Hybrid Syscall Execution Flow (ASCII)}]
User-space (Rust liburte)
    |
    v  syscall(URTE_xxx)
Kernel Entry
    |
    +--> [Guardrail Fast Path]  (Vector SIMD + small Rust functions)
    |         |
    |         v
    |    [TensorFlow Inference Path]  (PINN guardrail model / RL policy)
    |         |
    |         v
    +--> Result Aggregation (lock-free)
    |
    v  Return to user-space
\end{lstlisting}

\section{Hybrid Execution Layers}

\begin{enumerate}
    \item \textbf{Vector Layer (Rust + SIMD)}: Fast path for simple guardrail checks, resilience metric computation ($R = -\lambda_{\rm dom}$), and state vector operations using AVX-512 / NEON intrinsics via Rust `std::simd` or `packed_simd`.
    \item \textbf{TensorFlow Layer}: Complex guardrail evaluation, uncertainty quantification, and generative intervention design using TensorFlow Lite or full TensorFlow C API called from Rust via `tensorflow` crate or FFI.
    \item \textbf{Parallel Orchestration Layer}: Rust `rayon` or `tokio` + `crossbeam` for distributing work across multispectral thread pools.
\end{enumerate}

%==============================================================================
\chapter{Implementation with Rust Compiler}
%==============================================================================

\section{Rust as Primary Implementation Language}

\textbf{Advantages for URTE Kernel Threads:}
\begin{itemize}
    \item Memory safety without GC (critical for kernel/user-space shared components)
    \item Zero-cost abstractions for vector parallelism (`std::simd`, `rayon`)
    \item Excellent FFI to TensorFlow C API and POSIX
    \item Fearless concurrency for multispectral thread management
    \item Growing ecosystem for embedded/real-time (`embassy`, `tokio`, `crossbeam`)
\end{itemize}

\section{Example: Rust Implementation Skeleton for urte\_guardrail\_check}

\begin{lstlisting}[style=kernelschema, caption={Rust Implementation Outline (Hybrid)}]
use tensorflow::{Graph, Session, Tensor};
use packed_simd::f64x8; // Vector SIMD

#[repr(C)]
pub struct UrteGuardrailResult {
    pub decision: i32,
    pub reason_code: u32,
    pub required_actions: u64,
}

#[no_mangle]
pub extern "C" fn urte_guardrail_check(
    sv_fd: i32,
    req: *const UrteInterventionRequest,
    result: *mut UrteGuardrailResult,
) -> i32 {
    // 1. Fast vector path (Rust SIMD)
    let fast_path_score = unsafe { compute_vector_resilience(req) };

    if fast_path_score > THRESHOLD {
        // 2. TensorFlow path for complex evaluation
        let tf_result = run_tf_guardrail_model(req);
        unsafe { *result = tf_result; }
    } else {
        // Fast path result
    }
    0
}

fn compute_vector_resilience(req: *const UrteInterventionRequest) -> f64 {
    // AVX-512 / NEON accelerated computation
    let state: &[f64] = ...;
    let v = f64x8::from_slice_unaligned(state);
    // ... parallel reduction for lambda_dom, variance, etc.
    0.0
}
\end{lstlisting}

%==============================================================================
\chapter{Capella MBSE Modeling of Thread Environment}
%==============================================================================

\section{Arcadia Viewpoints for Threads}

\begin{itemize}
    \item \textbf{Logical Architecture}: Model thread pools as Logical Components with Ports for scale transitions and TRL Pull notifications.
    \item \textbf{Physical Architecture}: Allocate thread pools to Physical Nodes (Ground Control, Edge Devices, Satellites).
    \item \textbf{Functional Chains}: Model end-to-end syscall execution as Functional Chains with parallel branches (Vector + TensorFlow paths).
\end{itemize}

\section{Capella ASCII Representation Example}

\begin{lstlisting}[style=kernelschema, caption={Capella Logical Architecture - Thread Pools}]
Logical Architecture [LA] - URTE Thread Environment
+-- Logical Component: MultispectralThreadManager
|   Ports:
|     - ScaleTransitionPort
|     - TRLPullNotificationPort
|     - GuardrailEvaluationPort
|
+-- Logical Component: MolecularThreadPool (Scale 0-1)
|   Allocated Functions: Fast Vector Guardrail Checks
|
+-- Logical Component: EcosystemThreadPool (Scale 4-6)
|   Allocated Functions: Heavy TensorFlow PINN Inference
|
+-- Logical Component: HybridExecutionEngine
    Functions:
      - VectorSIMDPath
      - TensorFlowInferencePath
      - ResultAggregation
\end{lstlisting}

%==============================================================================
\chapter{Conclusion}
%==============================================================================

The URTE Kernel Threads Environment combines:
\begin{itemize}
    \item \textbf{Multispectral parallel thread management} across all regeneration scales
    \item \textbf{Hybrid TensorFlow + Vector (SIMD) computing} for high-performance system call execution
    \item \textbf{Rust} for safe, efficient, and parallel implementation
    \item \textbf{Capella/Arcadia} for formal, traceable MBSE modeling
\end{itemize}

This design preserves full POSIX compliance while delivering the performance and safety required for real-world deployment of the Universal Regeneration Therapy Ecosystem.

\vspace{0.8cm}
\begin{center}
\textit{--- End of URTE Kernel Threads Hybrid Parallel Specification v1.0 ---}
\end{center}

\end{document}
